\chapter{Teoria kategorii.}

W tym rozdziale wspiemy się na nowe poziomy abstrakcji.
Pojęcia takie jak grupa czy monoid powstały po to,
	by móc badać obiekty o podobnych cechach jednocześnie.
Zamiast udowadniać twierdzenia osobno dla wielomianów i liczb nautralnych,
	można udowodnić je dla monoidu---i to wystarczy.
W rozdziale \ref{ch:algebraic-structures} weszliśmy na wyższy poziom.
Zamiast definiować dla każdej struktury algebraicznej z osobna
	homomorfizmy i jądra czy udowadniać o nich twierdzenia,
	zrobiliśmy to jednocześnie używając algebry uniwersalnej.
Teoria kategorii jest bardzo podobnym pomysłem.
Tym razem jednak zamiast uogólniać pojęcia na jakąś rodzinę struktur (jak poprzednio na wszystkie algebry),
	uogólnimy pojęcia na całe dziedziny matematyki.

\section{Wstęp.}

Zaczniemy od pewnych podstawowych definicji i wprowadzenia do pojęcia kategorii.

\subsection{Klasy.}

Zaczniemy od pewnego pojęcia pierwotnego,
	które będzie nam potrzebne.
Pojęciem tym jest \textbf{klasa}.
Klasą nazywamy dowolną kolekcję obiektów.
Każdy zbiór tworzy klasę swoich elementów.
Klasy są jednak ogólniejsze niż zbiory.
Mówić będziemy m.in. o klasie wszystkich zbiorów,
	klasie wszystkich zbiorów nieskończonych czy skończonych,
	klasie wszyskich grup,
	klasie wszystkich przestrzeni liniowych...
W ogólności bpojęcia tego będziemy używać zamiast zbioru,
	kiedy zbiory będą zbyt małe.
Ale nam będzie potrzebny jakiś sposób ,,uchwycenia" wszystkich zbiorów.

Klasy dzielimy na dwa rodzaje: \textbf{zbiory} (które wszyscy wiedzą, czym są) i \textbf{klasy właściwe}.
Klasy właściwe to klasy,
	które nie mogą być zbiorami,
	np. klasa wszystkich zbiorów.
Jak wiadomo w końcu,
	zbiór wszystkich zbiorów nie istnieje.

\subsection{Kategorie.}

Definicja kategorii nie wymaga zbyt skomplikowanego aparatu.
Pojęcie klasy jest jedynym,
	które będzie nam tu potrzebne.
Przed przeczytaniem tej definicji należy wziąć jednak głęboki oddech.
Swoje życie matematyczne możemy podzielić na dwa etapy:
	przed zrozumieniem kategorii,
	oraz po ich zrozumieniu.
Powrotu już nie ma,
	a spojrzenie kategoryczne kategorycznie zmienia sposób myślenia o matematyce.

\begin{definition}
	\textbf{Kategorią} nazywamy parę klas $(O, A)$,
		gdzie $O$ jest klasą \textbf{obiektów}, a $A$ \textbf{strzałek}.
	Spełniają one następujące warunki
	\begin{itemize}
		\item Dla każdej strzałki $f$ istnieją obiekty $\dom(f)$ i $\cod(f)$
			nazywane odpowiednio \textbf{dziedziną} i \textbf{kodziedziną} strzałki,
			co zapisujemy
			\begin{equation}
				f: \dom(f) \to \cod(f).
			\end{equation}
		\item Dla każdych dwóch strzałek $f$, $g$, spełniających
			\begin{equation}
				\cod(f) = \dom(g)
			\end{equation}
			istnieje strzałka
			\begin{equation}
				g \circ f: \dom(f) \to \cod(g)
			\end{equation}
			nazywana \textbf{złożeniem} strzałek $f$ i $g$.
		\item Dla każdego obiektu $A$ istnieje strzałka
			\begin{equation}
				1_A: A \to A
			\end{equation}
			nazywana strzałką \textbf{identycznościową}.
		\item Dla dowolnych trzech strzałek
			\begin{align}
				f: A \to B,
				g: B \to C,
				h: C \to D,
			\end{align}
			spełniona jest własność łącznoci ze względu na składanie
			\begin{equation}
				h \circ (g \circ f) = (h \circ g) \circ f
			\end{equation}
		\item Dla dowolnej strzałki
			\begin{equation}
				f: A \to B
			\end{equation}
			spełnione jest
			\begin{equation}
				f \circ 1_A = f = 1_B \circ f.
			\end{equation}
	\end{itemize}
\end{definition}

Definicję zaczeliśmy od czegoś,
	co na myśl przywodzi graf skierowany.
Potem jednak zapis sugeruje,
	jakbyśmy działali na funkcjach oraz ich dziedzinach i kodziedzinach.
A kończymy z czymś podobnym do monoidu---mamy łączność i identyczność.
I rzeczywiście wszystkie te spojżenia są trafne.
Przejdźmy jednak najpierw do konkretnych przykładów kategorii.

\begin{example}
	Zbiory jako obiekty i funkcje jako strzałki tworzą kategorię.
	Rzeczywiście,
		wszystkie własności składania funkcji spełniają warunki,
		które spełniać muszą strzałki.
	Jako strzałkę identycznościową wystarczy wziąć funkcję identycznościową.
	Kategorię tę będziemy oznaczać przez \textbf{Set}.
\end{example}

\begin{example}
	Przestrzenie liniowe jako obiekty i przekształcenia liniowe jako strzałki tworzą kategorię.
	Udowadnia się tego identycznie, jak w przypadku zbiorów i funkcji
	(bo w końcu przestrzenie liniowe to też zbiory, a przekształcenia liniowe to też funkcje).
	Kategorię tę będziemy oznaczać przez \textbf{Lin}.
\end{example}

Nie w każdej kategorii jednak strzałki muszą odpowiadać funkcjom.

\begin{example}
	Zbiory jako obiekty i pary zbiorów $(A, B)$, że $A \subseteq B$ jako strzałki tworzą kategorię.
	Rzeczywiście,
		każde zawieranie
		\begin{equation}
			A \subseteq B
		\end{equation}
		możemy reprezentować przez strzałkę
		\begin{equation}
			(A, B): A \to B,
		\end{equation}
		a składanie strzałek i jego łączność jest możliwe
		ze względu na przechodniość relacji zawierania zbiorów.
		Strzałka identycznościowa istnieje,
			bo dla każdego zbioru $A$ mamy
		\begin{equation}
			A \subseteq A.
		\end{equation}
\end{example}

O ile w powyższych przykładach obiekty zawsze były jakimiś zbiorami,
	to wcale nie jest to konieczne.

\begin{example}
	Niech $~$ będzie dowolną relacją zwrotną i przechodnią na zbiorze $A$.
	Elementy $A$ jako obiekty i pary $(a, b)$, gdzie $a ~ b$,
		tworzą kategorię.
	Dowodzi się tego podobnie co faktu,
		że zbiory i ich zawieranie tworzą kategorię.
\end{example}

\begin{example}
	Klasy zdań matematycznych jako obiekty i dowody jako strzałki tworzą kategorię.
	Ograniczymi się do intuicyjnego wyjaśnienia.
	Każdy dowód ma pewne założenia i wnioski.
	Dowody można łączyć w ciągi dowodów,
		co jest oczywiście łączne.
	Dowód identycznościowy też jest dowodem ($t \implies t$).
\end{example}

\subsection{Diagramy.}

Tutaj będą kiedyś rysunki.

\subsection{Kilka podstawowych definicji.}

W tej sekcji zaczniemy powoli odkrywać,
	w jaki sposób kategorie pozwalają na uogólnianie definicji
	i twierdzeń na całe połacie matematyki.

\begin{definition}
	\textbf{Obiektem początkowym} kategorii nazywamy taki obiekt $0$,
		że dla każdego obiektu $A$ istnieje dokładnie jedna strzałka
		\begin{equation}
			\alpha_A: 0 \to A.
		\end{equation}
\end{definition}

\begin{example}
	W \textbf{Set} obiektem początkowym jest zbiór pusty.
	Rzeczywiście, dla każdego zbioru $A$ istnieje dokładnie jedna funkcja
	\begin{equation}
		\alpha_A: \emptyset \to A,
	\end{equation}
	która niczego nie mapuje.
\end{example}

\begin{definition}
	\textbf{Obiektem końcowym} kategorii nazywamy taki obiekt $1$,
		że dla każdego obiektu $B$ istnieje dokładnie jedna strzałka
		\begin{equation}
			\tau_A: A \to 1.
		\end{equation}
\end{definition}

\begin{example}
	W \textbf{Set} obiektem końcowym jest każdy zbiór jednoelementowy.
	Rzeczywiście, dla każdego zbioru $A$ i zbioru jednoelementowego $\{x\}$
	istnieje dokładnie jedna funkcja
	\begin{equation}
		\tau_A: A \to \{x\},
	\end{equation}
	którą możemy zdefiniować tylko na jeden sposób, t.j. dla każdego $a \in A$
	\begin{equation}
		\tau_A(a) = x.
	\end{equation}
\end{example}

\begin{definition}
	Strzałkę
	\begin{equation}
		f: A \to B
	\end{equation}
	nazywamy \textbf{mono} lub \textbf{monomorfizmem}, jeśli dla dowolnych strzałek
	\begin{equation}
		g, h: C \to A,
	\end{equation}
	że jeśli
	\begin{equation}
		fg = fh
	\end{equation}
	to
	\begin{equation}
		g = h.
	\end{equation}
\end{definition}

\begin{example}
	W \textbf{Set} monomorfizmy to dokładnie iniekcje.

	Najpierw pokażemy, że każdy monomorfizm to iniekcja.
	Niech $f: A \to B$ będzie dowolną funkcją mono.
	Niech $x, y \in B$.
	Załóżmy,
		że $f(x) = f(y)$.
	Chcemy pokazać,
		że $x = y$.

	Zdefiniujmy
	\begin{equation}
		\khi_x: \{0\} \to A, \quad \khi_x(0) = x.
		\khi_y: \{0\} \to A, \quad \khi_y(0) = y.
	\end{equation}

	Mamy treaz ...

	Ponieważ $f$ jest mono,
		to $\khi_x = \khi_y$.
	Więc $x = y$.

	Teraz pokażemy, że każda iniekcja to monomorfizm.
\end{example}

\begin{definition}
	Strzałkę
	\begin{equation}
		f: A \to B
	\end{equation}
	nazywamy \textbf{epi} lub \textbf{epimorfizmem}, jeśli dla dowolnych strzałek
	\begin{equation}
		g, h: B \to C,
	\end{equation}
	że jeśli
	\begin{equation}
		gf = hf
	\end{equation}
	to
	\begin{equation}
		g = h.
	\end{equation}
\end{definition}

\begin{example}
	W \textbf{Set} epimorfizmy to dokładnie surjekcje.
\end{example}

\subsection{Zadania.}

\section{Kategorie i język kategorii.}

\subsection{Zadania.}

\section{Różne konstrukcje kategoryczne.}

\subsection{Zadania.}

\section{Dualność i kilka innych pojęć.}

\subsection{Zadania.}

\input{03-category-theory/05-functors}
